\documentclass[../main.tex]{subfiles}
\begin{document}
\setchapterstyle{kao}
\setchapterpreamble[u]{\margintoc}
\setchapterimage[6.5cm]{Images/Flavours.png}
\chapter[Flavour Physics at Quark Level]{Flavour Physics at Quark Level\footnotemark[0]}
\labch{fpql}
\section{Symmetries at Quark Level}
In the quark model, states are bound states of quarks and anti-quarks:
\[
\hat{P}\ket{\bar{q}q',s=0,l=0}=\underset{\mathclap{\tikz \node {$\uparrow$} node [below=1ex] {\footnotesize  $\eta_q\eta_{q'}$};}}{(-1)}(-1)^l\ket{\bar{q}q',s=0,l=0}
\]
We have an intrinsic parity $(-)$ for $l=0$ and based just on parity conservation we have some processes which are not observed in nature:
\[
\left\{
\begin{aligned}
&0^-: \pi,k,\eta,\eta'\\
&1^-: \rho,k^*,\phi,\omega
\end{aligned}
\right.
\qquad
\left\{
\begin{aligned}
&\rho\xrightarrow[]{}\pi^+\pi^- \quad (-)\xrightarrow[]{}(-)^2(-)^1\\
&\rho\xrightarrow[]{}\pi^0\pi^0 \quad \text{ not observed in nature}
\end{aligned}
\right.
\]
We have seen in \refch{intro} how fields and operators behave under charge conjugation. If we apply this to the photons we get selection rules:
\[
\hat{C}\ket{\gamma}=-\ket{\gamma}\xrightarrow[]{}\hat{C}\ket{n\gamma}=(-)^n\ket{n\gamma}
\]
Hence, the transition $\gamma\gamma\xrightarrow[]{}3\gamma$ is forbidden (\href{https://en.wikipedia.org/wiki/Furry's_theorem}{Furry Theorem}).

We now bring our attention to the case of Positronium, which is a bound state of
\[
\ket{e^-,\Vec{p},\sigma;e^+,\Vec{p}',\sigma'}=\hat{a}^\dagger_\sigma(\Vec{p})\hat{b}^\dagger_{\sigma'}(\Vec{p}')\ket{0}
\]
Is it possible to treat Positronium as the Hydrogen atom? The masses are different: for Hydrogen, the relevant one is $m_e$ while for Positronium we have $m_e/2$. Nevertheless, this does not force us to treat them separately:
\[
E_n(H)=\frac{m_e\alpha^2}{2} \quad E_n(Ps)=\frac{\mu\alpha^2}{2}=\frac{m_e\alpha^2}{4}
\]
We work in the center of mass, hence $\Vec{p}'=-\vec{p}$. The setup is clear and we can apply charge conjugation:
\[
\hat{C}\ket{e^-,\Vec{p},\sigma;e^+,-\Vec{p},\sigma'}=\underbrace{\xi_{e^-}\xi_{e^+}}_{=1}\ket{e^+,\Vec{p},\sigma;e^-,-\Vec{p},\sigma'}=-\ket{e^-,-\Vec{p},\sigma';e^+,\Vec{p},\sigma}
\marginnote{We swapped momenta and spin of the two particles to get a factor $(-)$.}
\]
Applying charge conjugation to a general state in the form $\ket{n,l,m,s,s_z}$ gives us a factor $(-)^{l+s}$. Positronium can decay into photons:
\[
\left\{
\begin{aligned}
&n=1,l=0,s=0: (-)^{l+s}=+1\xrightarrow[]{}2\gamma, \tau=0.125\,\text{ns}\\
&n=1,l=1,s=1: (-)^{l+s}=-1\xrightarrow[]{}3\gamma, \tau=142.05\,\text{ns}
\end{aligned}
\right.
\marginnote{For this particular decay:
\[
\Gamma_{3\gamma}=\frac{e^2}{16\pi^2}\Gamma_{2\gamma}\sim10^{-4}\Gamma_{2\gamma}
\]}
\]
Moving back to our starting point, we observe how the decay channel of $\pi^0$ works: $\pi^0$ is a bound state of the form $\ket{q\bar{q},s=0,l=0}$ with a $(+)$ phase under charge conjugation. It then follows from this simple reasoning that the decay $\pi^0\xrightarrow[]{}3\gamma$ is forbidden (this has a branching ratio smaller than $10^{-8}$) while what we observe in nature is $\pi^0\xrightarrow[]{}2\gamma$.
\section{CP Violation in the Standard Model}
For parity, charge conjugation and CP application, we have to keep in mind the following table:
\begin{table}[h]
    \centering
    {\renewcommand{\arraystretch}{1}
    \begin{tabular}{c|c|c|c}
    Fields & $\hat{P}(x\xrightarrow[]{}\pazocal{P}x)$ & $\hat{C}$ & $\hat{CP}(x\xrightarrow[]{}\pazocal{P}x)$\\
    \hline
    $\bar{\Psi}\chi$ & $\bar{\Psi}\chi$ & $\bar{\chi}\Psi$ & $\bar{\chi}\Psi$\\
    $\bar{\Psi}_R\chi_L$ & $\bar{\Psi}_L\chi_R$ & $\bar{\chi}_R\Psi_L$ & $\bar{\chi}_L\Psi_R$\\
    $\bar{\Psi}_L\chi_R$ & $\bar{\Psi}_R\chi_L$ & $\bar{\chi}_L\Psi_R$ & $\bar{\chi}_R\Psi_L$\\
    $\bar{\Psi}\gamma^\mu\chi$ & $\bar{\Psi}\gamma_\mu\chi$ & $-\bar{\chi}\gamma^\mu\Psi$ & $-\bar{\chi}\gamma_\mu\Psi$ \\
    $\bar{\Psi}_L\gamma^\mu\chi_L$ & $\bar{\Psi}_R\gamma_\mu\chi_R$ & $-\bar{\chi}_R\gamma^\mu\Psi_R$ & $-\bar{\chi}_L\gamma_\mu\Psi_L$\\
    $\bar{\Psi}_R\gamma^\mu\chi_R$ & $\bar{\Psi}_L\gamma_\mu\chi_L$ & $-\bar{\chi}_L\gamma^\mu\Psi_L$ & $-\bar{\chi}_R\gamma_\mu\Psi_R$\\
    \hline
    \end{tabular}
    }
    \caption*{}
    \label{tab:my_label}
\end{table}\\
For every term there is a phase factor $e^{i(\phi_\chi-\phi_\Psi)}$ associated,\\
e.g. $\bar{\Psi}\chi(x)\xrightarrow[]{\hat{CP}}e^{i(\phi_\chi-\phi_\Psi)}\bar{\chi}\Psi(\pazocal{P}x)$. Now we apply $\hat{CP}$ to all the stuff we already know.
\[
\left\{
\begin{aligned}
&\hat{CP}W_\mu^+(x)\hat{CP}^\dagger=-e^{+i\xi_W}\pazocal{P}_\mu^\nu W_\mu^-(\pazocal{P}x) \quad &&\hat{CP}W_\mu^-(x)\hat{CP}^\dagger=-e^{-i\xi_W}\pazocal{P}_\mu^\nu W_\mu^+(\pazocal{P}x)\\
&\hat{CP}Z_\mu(x)\hat{CP}^\dagger=-\eta_Z\pazocal{P}_\mu^\nu Z_\nu(\pazocal{P}x) \quad && \hat{CP}A_\mu(x)\hat{CP}^\dagger=-\pazocal{P}_\mu^\nu A_\nu(\pazocal{P}x)\\
&\hat{CP}G_\mu^a(x)\hat{CP}^\dagger=-\eta_a\pazocal{P}_\mu^\nu G_\nu^a(\pazocal{P}x) \quad && \hat{CP}h(x)\hat{CP}^\dagger=\eta_h h(\pazocal{P}x)\\
&\hat{CP}u_i\hat{CP}^\dagger=e^{i\phi_{u_i}}\gamma^0C\bar{u}_i^T \quad && \hat{CP}d_i\hat{CP}^\dagger=e^{i\phi_{d_i}}\gamma^0C\bar{d}_i^T
\end{aligned}
\right.
\]
Is the SM \textbf{invariant} under $\hat{CP}$ or not?  Is it possible to find $\xi_W, \eta_Z, \eta_h, \phi_{u_i}, \phi_{d_i}$ such that $\hat{CP}\pazocal{L}(x)\hat{CP}^\dagger=\pazocal{L}(\pazocal{P}x)?$ The Standard Model Lagrangian can be decomposed as $\pazocal{L}_{SM}=\pazocal{L}_{kin}+\pazocal{L}_{Higgs}+\pazocal{L}_{Yuk}$. 

The typical term containing the $Z$ boson is in the form
\[
\sum_{f=u,d;A=L,R}\bar{f}_A\gamma^\mu f_AZ_\mu g_A^f\xleftarrow[]{}\text{real and constant coefficients}
\]
We apply $\hat{CP}$ to it: fermions appear as bilinears, the phase is cancelled because there is $\bar{f}_Af_A$ so we just focus on the sign. Under $\hat{CP}$ we get $\bar{f}_A\gamma^\mu f_A\xrightarrow[]{}(-\bar{f}_A\gamma^\mu f_A)e^{i\phi_{f_A}}e^{-i\phi_{f_A}}$. For the $Z$ we have $-(Z^\mu)\eta_Z$, therefore it has to be $\eta_Z=+1$. 

For the Higgs we have two pieces: $\pazocal{L}_{Higgs}=-m_h^2h^2+avh^3+\lambda h^4$ and the term $\sim h^3$ tells us that it has to be $\eta_h=+1$. 

The Yukawa sector does not give us a relevant contribution:
\[
\hat{CP}\pazocal{L}_{Yuk}\hat{CP}^\dagger=\hat{CP}Y_fh\bar{f}f\hat{CP}^\dagger=\hat{CP}Y_fh\hat{CP}^\dagger\hat{CP}\bar{f}f\hat{CP}^\dagger
\]
We apply $\hat{CP}$ to the $h$ and get $\eta_h=+1$ as seen before, while $\bar{f}f$ gives a +1 sign so everything is fine.

We are now left with $\xi_W, \phi_{u_i}$ and $\phi_{d_i}$:
\[
\hat{CP}\left[\frac{g}{\sqrt{2}}W_\mu^+\bar{u}_L^i\gamma^\mu d^j_LV_{ij}+\frac{g}{\sqrt{2}}W_\mu^-\bar{d}_L^j\gamma^\mu u_L^iV_{ji}^*\right]\hat{CP}^\dagger
\]
We insert $\hat{CP}^\dagger\hat{CP}=1$ in the expression above and what we find is:
\begin{align*}
&\frac{g}{\sqrt{2}}V_{ij}\left[-e^{+i\xi_W}(W^-)^\mu e^{i(\phi_{d_j}-\phi_{u_i})}\left(-\bar{d}_L^j\gamma_\mu u_L^i\right)\right]\\
+&\frac{g}{\sqrt{2}}V_{ji}^*\left[-e^{-i\xi_W}(W^+)^\mu e^{i(\phi_{u_i}-\phi_{d_j})}\left(-\bar{u}_L^i\gamma_\mu d_L^j\right)\right]
\end{align*}
In order to get CP invariance, it has to be $V_{ij}^*=V_{ij}e^{i(\xi_W+\phi_{d_j}-\phi_{u_i})}$: can we fix these phases such that this combination becomes real?\\
The matrix $V$$\in$U$(3)$ and we have 9 parameters: 3 angles+6 phases. However, in our case we have only 7 parameters: $\overset{(1)}{\xi_W}, \overset{(3)}{\phi_{u_i}}, \overset{(3)}{\phi_{d_j}}$ but some of them are useless. In principle, we have infinite solutions corresponding to 2 unbroken symmetries: U$(1)_B$, in which $\phi_u=\phi_d$ and we cannot use that because there are 5 phases; U$(1)_Q$, corresponding to electromagnetism with the phases in the ratio +1, -2/3, -1/3. 

We want to find CP violation in the unbroken phase, so with massless fields. For $\pazocal{L}_{kin}+\pazocal{L}_{Higgs}$ we have a U$(3)^3$ symmetry: U$(3)_{Q_L}\times$U$(3)_{u_R}\times$U$(3)_{d_R}$. This symmetry is broken by adding $\pazocal{L}_{Yuk}$, but before adding this piece we define a new CP transformation, which includes the U$(3)^3$ operators, and ask whether $\pazocal{L}_{Yuk}$ is invariant under the generalized CP transformation defined as follow:
\[
\left\{
\begin{aligned}
&\hat{CP}Q_L^i\hat{CP}^\dagger=(K_L)^i_je^{i\phi_{Q_L^i}}\gamma^0C(\bar{Q}_L^i)^T\\
&\hat{CP}u_R^i\hat{CP}^\dagger=(K_u)^i_je^{i\phi_{u_R^i}}\gamma^0C(\bar{u}_R^i)^T\\
&\hat{CP}d_R^i\hat{CP}^\dagger=(K_d)^i_je^{i\phi_{d_R^i}}\gamma^0C(\bar{d}_R^i)^T
\end{aligned}
\right.
\quad 
K_L,K_u,K_d\in\text{U}(3)
\]
We can now consider $-\Tilde{H}\bar{Q}_L^i(Y_u)_{ij}u_R^j-\Tilde{H}\bar{Q}_L^i(Y_d)_{ij}d_R^j$ and focus on flavour indices: under CP we have that
\begin{equation}
\labeq{p1}
\left\{
\begin{aligned}
&Y_u\xrightarrow[]{}K_L^\dagger Y_uK_u\stackrel{?}{=}Y_u^*\\
&Y_d\xrightarrow[]{}K_L^\dagger Y_dK_d\stackrel{?}{=}Y_d^*
\end{aligned}
\right.
\end{equation}
So far, we rephrased the question about CP invariance in: are there $K_L,K_u$ and $K_d$ such that the condition above holds true? This is the same as:
\begin{equation}
\labeq{p2}
\left\{
\begin{aligned}
&K_L^\dagger Y_uY_u^\dagger K_L=(Y_uY_u^\dagger)^*\\
&K_L^\dagger Y_dY_d^\dagger K_L=(Y_dY_d^\dagger)^*
\end{aligned}
\right.
\end{equation}
In fact, by using \refeq{p1} we can show that we get \refeq{p2}:
\[
K_L^\dagger Y_uY_u^\dagger K_L=\underbrace{K_L^\dagger Y_u{\color{red}K_u}}_{Y_u^*}\underbrace{{\color{red}K_u^\dagger}Y_u^\dagger K_L}_{(Y_u^*)^\dagger}=Y_u^*(Y_u^*)^\dagger=(Y_uY_u^\dagger)^* \quad \checkmark
\]
We just need to prove that $K_u=Y_u^{-1}K_LY_u^*$ is unitary:
\[
K_uK_u^\dagger=Y_u^{-1}K_L\underbrace{Y_u^*Y_u^T}_{=(Y_uY_u^\dagger)^*}K_L^\dagger(Y_u^{-1})^T=Y_u^{-1}K_LK_L^\dagger Y_uY_u^\dagger K_LK_L^\dagger(Y_u^{-1})^T=\mathbb{1} \quad \checkmark
\]
If CP is conserved, then \refeq{p2} holds true. Moreover, if CP is conserved we have that $\Tr{[H_u,H_d]^r}=0$ when $r$ is odd, where we have defined $H_u:=Y_uY_u^\dagger$ and $H_d:=Y_dY_d^\dagger$, which are hermitian. To prove the second statement, we compute:
\begin{align*}
K_L^\dagger[H_u,H_d]K_L&=K_L^\dagger H_uH_dK_L-K_L^\dagger H_dH_uK_L\\
&=(K_L^\dagger H_uK_L)(K_L^\dagger H_dK_L)-(K_L^\dagger H_dK_L)(K_L^\dagger H_uK_L)\\
&\underset{\mathclap{\tikz \node {$\uparrow$} node [below=1ex] {\footnotesize \refeq{p2} };}}{=}H_u^*H_d^*-H_d^*H_u^*=[H_u^T,H_d^T]=-[H_u,H_d]^T
\end{align*}
It follows that $\Tr{[H_u,H_d]}=-\Tr{[H_u,H_d]}=0$. This was the case for $r=1$, now we generalize to any value of $r$:
\begin{align*}
\left(K_L^\dagger[H_u,H_d]K_L\right)^r&=K_L^\dagger[H_u,H_d]K_LK_L^\dagger[H_u,H_d]K_L\dots K_L^\dagger[H_u,H_d]K_L\\
&=K_L^\dagger[H_u,H_d]^rK_L=(-1)^r\left([H_u,H_d]^T\right)^r
\end{align*}
It then follows that $\Tr{K_L^\dagger[H_u,H_d]^rK_L}=\Tr{[H_u,H_d]^r}=(-1)^r\Tr{[H_u,H_d]^r}$. This tells us that if we have a $r$ odd such that $\Tr{[H_u,H_d]^r}\neq0$, CP is violated. For $N_F=3$, we have that CP violation$\Rightarrow\Tr{[H_u,H_d]^3}=0$, while in all the other cases we just have $\Tr{[H_u,H_d]^3}\neq0\Rightarrow$CP violation. We can explicitly compute it for $r=3, N_F=3$:
\[
\Tr{[H_u,H_d]^3}=6i(y_t^2-y_c^2)(y_t^2-y_u^2)(y_c^2-y_u^2)(y_b^2-y_s^2)(y_b^2-y_d^2)(y_s^2-y_d^2)J
\]
$J$ is the \href{https://en.wikipedia.org/wiki/Cecilia_Jarlskog}{Jarlskog} invariant
\marginnote{When pronouncing \textit{Jarlskog}, you might accidentally summon a demon in the extra-dimension (bad thing).
\includegraphics[]{Images/jarlskog.jpeg}
From \href{https://cerncourier.com/a/viewpoint-theory-at-cern-turns-62/}{Cern Courier}: Cecilia Jarlskog with colleagues at the Nordic Institute of Theoretical Physics (NORDITA) in Copenhagen, in the early 1980s.
}. The global symmetry G=U$(3)_{Q_L}\times$U$(3)_{u_R}\times$U$(3)_{d_R}$ given by $\pazocal{L}_{kin}+\pazocal{L}_{Higgs}$ gets broken into H=U$(1)_B$ when we add $\pazocal{L}_{Yuk}$. This introduces $2N^2+2N^2=18+18=36$ naive parameters, corresponding to $Y_u$ and $Y_d$. For G, we have 27 generators, while for H we have just 1: this means a total number of broken generators equal to 26. Finally, we know that
\[
\#\text{physical parameters}&=\#\text{naive}-\#\text{broken}=36-26=10=6+3+1
\]
This corresponds to 6 masses of quarks+3 mixing angles+CP violation of the CKM.
\section{Parametrization of the CKM Matrix}
Most generally, the CKM matrix can be written as:
\[
V_{CKM}=\left(\begin{array}{ccc}
    V_{ud} & V_{us} & V_{ub} \\
    V_{cd} & V_{cs} & V_{cb} \\
    V_{td} & V_{ts} & V_{tb}
\end{array}\right)
\]
This is unitary with 3 mixing angles ($\theta_{12},\theta_{23},\theta_{13}$)+1 phase ($\delta$)\marginnote{Actually there are 5 phases but 4 of them are unphysical.} and therefore it is possible to express it in the following way:
\[
V_{CKM}=R_{23}(\theta_{23})\mqty(\dmat{e^{-i\delta},1,1})R_{13}(\theta_{13})\mqty(\dmat{e^{+i\delta},1,1})R_{12}(\theta_{12})
\]
The three angles and the phase are such that, according to \href{https://it.wikipedia.org/wiki/Nicola_Cabibbo}{Cabibbo}, we have:
\[
\left\{
\begin{aligned}
&\sin\theta_{12}=0.225\pm0.001:=\lambda \quad && \sin\theta_{23}=0.04200\pm0.00059:=\lambda^2\\
&\sin\theta_{13}=0.00368\pm0.00010 \quad && \delta[^\circ]=66.9\pm2.0\xrightarrow[]{}\sin\theta_{13}e^{i\delta}:=A\lambda^3(\rho+i\eta)
\end{aligned}
\right.
\]
It is now possible to express the CKM matrix in function of the new 4 parameters just defined $(A,\lambda,\rho,\eta)$, with $\lambda$ being small and used as in a sort of perturbative expansion.\marginnote{The subscript W is due to the fact that the parametrization had been done by \href{https://en.wikipedia.org/wiki/Lincoln_Wolfenstein}{Wolfenstein}.

\includegraphics[]{Images/Wolfenstein.jpeg}
From \href{https://www.post-gazette.com/news/obituaries/2015/04/04/Obituary-Lincoln-Wolfenstein-Internationally-known-physicist-who-taught-at-CMU/stories/201504040050}{Pittsburgh Post-Gazzette}: Lincoln Wolfenstein (February 10, 1923, Cleveland, Ohio – March 27, 2015, Oakland, California) was an American particle physicist who studied the weak interaction. 
}
\begin{align*}
V_W&=\left(\begin{array}{ccc}
    1-\lambda^2/2 & \lambda & A\lambda^3(\rho-i\eta) \\
    -\lambda & 1-\lambda^2/2 & A\lambda^2 \\
    A\lambda^3(1-\rho-i\eta) & -A\lambda^2 & 1
\end{array}\right)
+\pazocal{O}(\lambda^4)\\
V_W&\sim
\left(\begin{array}{ccc}
    1 & \lambda & \lambda^3 \\
    \lambda & 1 & \lambda^2 \\
    \lambda^3 & \lambda^2 & 1
\end{array}\right)
\end{align*}
The new parameters have been calculated and they turn out to be:
\[
\left\{
\begin{aligned}
&\lambda=0.22574\pm0.00084 \quad &&\rho=0.152\pm0.014\\
&A=0.826\pm0.012 \quad &&\eta=0.357\pm0.010
\end{aligned}
\right.
\]
$V$ is unitary, therefore we have that $\sum_{i=1}^3V_{id}V_{is}^*=0$ and we can construct 6 relations of this type using rows and columns. In this case, we have $V_{ud}V_{us}^*\sim\lambda$, $V_{cd}V_{cs}^*\sim\lambda$ and $V_{td}V_{ts}^*\sim\lambda^5$. We can represent this equality in the complex plane obtaining a unitarity triangle.
\marginnote{
\includegraphics{Images/uni_triangle_1.pdf}
Representation of the unitarity triangle.
}
We write the 6 possible relations using all the rows and columns in the matrix, obtaining 6 different triangles which all have the same area equal to $J/2$. With this parametrization, $J$ turns out to be equal to:
\[
J=\cos\theta_{12}\cos\theta_{23}\cos\theta_{13}^2\sin\theta_{12}\sin\theta_{23}\sin\theta_{13}\sin\delta\simeq\lambda^6A^2\eta
\]
In the example we made, the sides of the triangles are of different scaling, a better choice instead is to use $\sum_{i=1}^3V_{id}V_{ib}^*$. In this way, we obtain $\{V_{ud}V_{ub}^*,V_{cd}V_{cb}^*,V_{td}V_{tb}^*\}\sim\lambda^3$ and in order to have a better representation on the complex plane we can redefine everything as:
\[
\frac{V_{ud}V_{ub}^*}{V_{cd}V_{cb}^*}+1+\frac{V_{td}V_{tb}^*}{V_{cd}V_{cb}^*}=0
\]
\begin{figure}[h]
    \centering
    \includegraphics[width=0.8\textwidth]{Images/uni_triangle_2.pdf}
    \caption{Unitarity triangle after the redefinition. The notation may vary: in California there is $(\gamma, \alpha, \beta)$ while in Japan they call them $(\phi_3,\phi_2,\phi_1)$. We introduced $R_u=\left|\frac{V_{ud}V_{ub}^*}{V_{cd}V_{cb}^*}\right|$ and $R_t=\left|\frac{V_{td}V_{tb}^*}{V_{cd}V_{cb}^*}\right|$}
    \label{fig:my_label}
\end{figure}\\
Experimentally, what are measured are the angles $(\alpha,\beta,\gamma)$ in order to obtain $(\rho,\eta)$. Moreover, the measures give us a check of consistency: if we find a strip of measurements out of the expected zone, we find an inconsistency which can be due to new physics. CP violation is present in the Standard Model but there are cosmological reasons which tell us there must be some other source of CP violation in the Universe.
\section{Flavour Changing Neutral Currents (FCNC)}
\subsection{FCNC at Tree Level}
We have chosen a basis for the quarks where the off-diagonal entries come from the $W$ boson, which we know it is charged. Therefore, the neutral boson interactions are flavour conserving. Let's look at this from the experimental side and suppose we know nothing about the Standard Model. We just know experimental data, which tell us:
\begin{table}[h]
    \centering
    \begin{tabular}{c|c}
    Flavour Changing Charged Current & Flavour Changing Neutral Current \\
    \hline\hline
    BR$(K^+\xrightarrow[]{}\mu^+\nu_\mu)=64\%$ & BR $(K_L^0\xrightarrow[]{}\mu^+\mu^-)=7\times10^{-3}$\\
    $|\Delta S|=1$ & $|\Delta S|=1$\\
    $|\Delta Q|=1$ & $|\Delta Q|=0$\\
    \hline
    BR$(B^-\xrightarrow[]{}D^0l\nu)=2.3\%$ & BR $(B^-\xrightarrow[]{}K^*^-l^+l^-)=5\times10^{-7}$\\
    $|\Delta B|=1$ & $|\Delta B|=1$\\
    $|\Delta Q|=1$ & $|\Delta Q|=0$\\
    \hline
    \end{tabular}
    \caption*{}
    \label{tab:my_label}
\end{table}
\newline
Why is there no tree-level neutral current in the SM?\\
The possible tree-level neutral mediators of the SM are photons $\gamma$, gluons $G$, the Higgs $h$ and the $Z$ boson. Photons and gluons cannot give FCNC at tree-level because of gauge invariance: as long as QCD and QED are unbroken, all the up and the down quarks will have the same gauge representations, forcing the kinetic terms to be universal.
\paragraph{Higgs FCNC} The Higgs $h$ is "accidentally" alone: how can it be flavour conserving? It couples to fermions via the Yukawa term, take for example\\ $(Y^u)_{ij}\Tilde{H}\bar{Q}_L^iu^j+\text{h.c.}$ and expand $H$ around its v.e.v., which in the unitary gauge means to consider:
\[
H=\left(\begin{array}{c}
    \frac{v+h}{\sqrt{2}} \\
    0
\end{array}\right)
\quad
\Tilde{H}=\left(\begin{array}{c}
    0\\
    \frac{v+h}{\sqrt{2}}
\end{array}\right)
\]
For the mass term we have $(m^u)_{ij}=(Y^u)_{ij}v/\sqrt{2}$, while for the coupling we obtain $(Y_u)_{ij}=(Y^u)_{ij}/\sqrt{2}$. Now we move to a basis in which the mass is diagonal:
\[
\left\{
\begin{aligned}
&u_L=V_L^uu_L'\\
&u_R=V_R^uu_R'
\end{aligned}
\right.
\Rightarrow
\left\{
\begin{aligned}
&\hat{m}^u=V_L^umV_R^u\\
&\hat{Y}^u=V_L^uY_u(V_R^u)^\dagger
\end{aligned}
\right.
\]
The absence of FCNC at tree-level is due to the fact that we have only one Higgs and we are in 4 dimensions. When looking for new physics we can either add non-renormalizable operators or add another Higgs. We now proceed following the first idea, taking the usual Lagrangian and adding an extra-term:
\[
C_{ij}\frac{\Tilde{H}Q_L^iu_R^jH^\dagger H}{\Lambda^2}+(Y^u)_{ij}\Tilde{H}\bar{Q}_L^iu^j+\text{h.c.}
\]
This modifies both the mass and the coupling:
\[
\left\{
\begin{aligned}
&m_u=Y^u\frac{v}{\sqrt{2}}+C\left(\frac{v}{\sqrt{2}}\right)^3\frac{1}{\Lambda^2}\\
&Y_u=\frac{Y^u}{\sqrt{2}}+C\left(\frac{v}{\sqrt{2}}\right)^2\frac{1}{\Lambda^2}3\xleftarrow[]{}\text{possible combinations}
\end{aligned}
\right.
\]
If we try to diagonalize it as before, for the mass we obtain the same result while for the interaction we get:
\[
Y_h=V_L^u\left(\frac{Y^u}{\sqrt{2}}+C\frac{v^2}{2\Lambda^2}\right)(V_R^u)^\dagger+2V_L^uC\frac{v^2}{2\Lambda^2}(V_R^u)^\dagger
\]
This is not diagonal: non-renormalizable operators can give rise to FCNC at tree-level. $C_{ij}\xrightarrow[]{}(Y^u)_{ij}$ apparently cures the problem of flavour violation since $m_u\propto Y_u$ and $Y_h\propto Y_u$.\\
Now we explore the possibility of adding a second Higgs to the SM and see what happens. \marginnote{Models with two Higgs fields are important since the minimal supersymmetric models require at least two Higgs fields.}
\[
\left\{
\begin{aligned}
&H_1\sim(1,2,1/2) \quad &&\Tilde{H}_1=i\sigma_2H_1^*\\
&H_2\sim(1,2,1/2) \quad &&\Tilde{H}_2=i\sigma_2H_2^*
\end{aligned}
\right.
\]
The potential $V(H)$ gets modified into $V(H_1,H_2)$, with the most general v.e.v. pattern give by:
\[
\langle H_1 \rangle=\left(\begin{array}{c}
     0 \\
     v_1/\sqrt{2}
\end{array}\right)
\quad
\langle H_2 \rangle=\left(\begin{array}{c}
     W^+ \\
     v_2/\sqrt{2}
\end{array}\right)
\]
The parameters $v_1, v_2$ and $W^+$ are function of the potential $V$. Since it is designed to allow U$(1)_Q$ symmetry, then it must be $W^+=0$. We can now count the degrees of freedom of a model with two complex SU(2) doublets, in an attempt to understand what are the fields which arise. Two complex doublets gives 4+4=8 degrees of freedom. The initial symmetry SU$(2)_L\times$U$(1)_Y$ is broken into U$(1)_Q$: this means that we have $2^2-1=3$ broken generators, corresponding to 3 NGB.
\[
\text{\# degrees of freedom}=8=\underset{\mathclap{\tikz \node {$\uparrow$} node [below=1ex] {\footnotesize $H_1$ };}}{4}+\overset{\mathclap{\tikz \node {$\downarrow$} node [above=1ex] {\footnotesize $H_2$};}}{4}=\underset{\mathclap{\tikz \node {$\uparrow$} node [below=1ex] {\footnotesize  NGB};}}{3}+5=3+\underset{\mathclap{\tikz \node {$\uparrow$} node [below=1ex] {\footnotesize charged Higgs };}}{2}+\overset{\mathclap{\tikz \node {$\downarrow$} node [above=1ex] {\footnotesize neutral Higgs };}}3
\]
It is possible to assign to the neutral Higgs a CP parity, with 2 of them being even and 1 odd. The usual Higgs boson $h$ is the lightest CP even state. The generic Higgs field is given by:
\[
H_i=\left(\begin{array}{c}
    \phi_i^+\\
    \phi_i
\end{array}\right)
\underset{\mathclap{\tikz \node {$\uparrow$} node [below=1ex] {\footnotesize expand it around the v.e.v. };}}{\simeq}
\left(\begin{array}{c}
    \phi_i^+\\
    \frac{v_i+\rho_i+i\eta_i}{\sqrt{2}}
\end{array}\right)
\quad
i=1,2
\]
Guided by charge and CP parity, we can find the remaining 5 degrees of freedom.
\[
\left\{
\begin{aligned}
&(\phi_1^+,\phi_2^+)\xrightarrow[]{}(G^\pm,H^\pm) \quad &&G^\pm\,\text{: 2 NGB eaten by $W_L^\pm$},\, H^\pm\,\text{: 2 charged Higgs fields}\\
&\eta_i\xrightarrow[]{}(G^0,A)\quad &&G^0\,\text{: NGB eaten by $Z^0_L$},\, A\,\text{: neutral, CP odd}\\
&\rho_i\xrightarrow[]{}(h,H) \quad &&\text{both neutral and CP even. $h$: usual Higgs field, $H$: new heavy Higgs field}
\end{aligned}
\right.
\]
$\phi_i^+$ has charge equals to $\pm1$, hence we have terms like $m_{11}\phi_1^+\phi_1+m_{12}\phi_1^+\phi_2$ but not something like $\phi_1^+\rho$: we can only mix particles with the same symmetry. Now we bring our attention to the two new Higgs, $h$ and $H$:
\[
\left\{
\begin{aligned}
&\rho_1=\cos\alpha h+\sin\alpha H\\
&\rho_2=-\sin\alpha h+\cos\alpha H
\end{aligned}
\right.
\]
The parameter $\alpha$ depends on the structure of $V(H_1,H_2)$.

We can prove the existence of FCNC with 2 Higgs fields by doing the following exercise:
\begin{itemize}
    \item Write the most general $\pazocal{L}_{Yuk}$ including $H_1$ and $H_2$:
    \begin{align*}
    &-(Y_u^1)_{ij}\bar{Q}_L^iu_R^j\Tilde{H}_1-(Y_u^2)_{ij}\bar{Q}_L^iu_R^j\Tilde{H}_2+\text{h.c.}\\
    &-(Y_d^1)_{ij}\bar{Q}_L^id_R^jH_1-(Y_d^2)_{ij}\bar{Q}_L^id_R^jH_2+\text{h.c.}
    \end{align*}
    \item Move to the unitary gauge:
    \[
    H_i=\left(\begin{array}{c}
    \phi_i^+\\
    \frac{v_i+\rho_i+i\eta_i}{\sqrt{2}}
    \end{array}\right)
    \xrightarrow[]{}
    \left(\begin{array}{c}
    0\\
    \frac{v_i+\rho_i}{\sqrt{2}}
    \end{array}\right)
    \quad
    \begin{cases}
    \text{$v_i$ mass term}\\
    \text{$\rho_i$ proportional to the Higgs (we do not consider the heavy field $H$)}
    \end{cases}
    \]
    \item Compute the mass matrix and the couplings:
    \[
    \left\{
    \begin{aligned}
    &(m_u)_{ij}=\frac{1}{\sqrt{2}}\left[(Y_u^1)_{ij}v_1+(Y_u^2)_{ij}v_2\right]\\
    &(Y_h^u)_{ij}=\frac{1}{\sqrt{2}}\left[(Y_u^1)_{ij}\cos\alpha+(Y_u^2)_{ij}(-\sin\alpha)\right]
    \end{aligned}
    \right.
    \]
\end{itemize}
By introducing $u_L=V_Lu_L'$ and $u_R=V_Ru_R'$:
\[
\pazocal{L}_{Yuk}\supset-\frac{1}{\sqrt{2}}u_L'(V_L^\dagger m_uV_R)u_R'-\frac{1}{\sqrt{2}}u_L'h(V_L^\dagger Y_hV_R)u_R'
\]
The mass term gets diagonalized, however $Y_h$ does not get diagonalized and we can have interactions like:
\begin{figure}[h]
    \centering
    \includegraphics[width=0.5\textwidth]{Images/feynman-1.pdf}
    \caption{Type of interaction we can have in this model.}
    \label{fig:my_label}
\end{figure}\\
The 2 Higgs model is important since we have two Higgs also in the SUSY model: one Higgs couples only to up quarks $(\Tilde{H}_1\to H_u)$ and one couples only to down quarks $(H_2\to H_d)$.
\[
\pazocal{L}_{Yuk}=-(Y_u)_{ij}\bar{Q}_L^iu_R^jH_u+\text{h.c}-(Y_d)_{ij}\bar{Q}_L^id_R^jH_d+\text{h.c.}
\]
With this assumption, we removed two couplings. Here interactions and mass re-align, giving us no flavour violation. The realignment of mass and interaction is always present in the SM with 1 Higgs.
\paragraph{$Z$ boson FCNC}
Let's look at the following toy model:
\begin{enumerate}
    \item Group symmetry: G=G$_{SM}=$U$(1)_{Q_L}\times$U$(1)_{s_L}\times$U$(1)_{u_R}\times$U$(2)$
    \item Field content: 
    \[
    \left\{
    \begin{aligned}
    &Q_L\sim(3,2,+1/6) \quad &&s_L\sim(3,1,-1/3)\\
    &u_R\sim(3,1,+2/3) \quad &&s_R,d_R\sim(3,1,-1/3)\\
    &H\sim(1,2,+1/2)
    \end{aligned}
    \right.
    \]
    \item Most general Lagrangian in the Yukawa sector (for the other sectors it is just the covariant derivative):
    \[
    \pazocal{L}_{Yuk}=-\left[Y_u\bar{Q}_L\Tilde{H}u_R+Y_d\bar{Q}_Ld_RH+Y_s\bar{Q}_L\Tilde{H}s_R+m_s\bar{s}_Ls_R+m_d\bar{s}_Ld_r+\text{h.c.}\right]
    \]
\end{enumerate}
G has dimension 7=1+6 and gets broken into H$=U(1)_B$ which has dimension 1. We have 10 naive parameters, corresponding to 5 moduli+5 phases:
\[
\text{\#physical parameters}=\text{\#naive}-\text{\#broken}=10-6=\underset{\mathclap{\tikz \node {$\uparrow$} node [below=1ex] {\footnotesize moduli };}}{5}+\overset{\mathclap{\tikz \node {$\downarrow$} node [above=1ex] {\footnotesize  phases};}}{5}-\underset{\mathclap{\tikz \node {$\uparrow$} node [below=1ex] {\footnotesize modulus};}}{1}-\overset{\mathclap{\tikz \node {$\downarrow$} node [above=1ex] {\footnotesize phases };}}{5}=4
\]
BY removing all the phases, we obtained CP invariance. With a rotation inside U(2) we remove another parameter: $\bar{s}_L(m_ss_R+m_dd_R)=m_s\bar{s}_Ls_R$. We are now interested in diagonalizing the mass term:\marginnote{In the matrix $M$ there should have been $m_d$ instead of 0 but we killed it with the U(2) rotation.}
\[
\pazocal{L}_{mass}=-\left[\underset{\mathclap{\tikz \node {$\uparrow$} node [below=1ex] {\footnotesize  $\frac{Y_uv}{\sqrt{2}}$};}}{m_u}\bar{u}_Lu_R+\left(\begin{array}{cc}
    \bar{d}_L & \bar{s}_L
\end{array}\right)
\underbrace{\left(\begin{array}{cc}
    \frac{Y_dv}{\sqrt{2}} & \frac{Y_sv}{\sqrt{2}} \\
    0 & m_s
\end{array}\right)}_{M}
\left(\begin{array}{c}
    d_R\\
    s_R
\end{array}\right)\right]
\]
We can diagonalize the mass matrix through:
\[
\left(\begin{array}{c}
     d_R\\
     s_R
\end{array}\right)=V_R^\dagger\left(\begin{array}{c}
     d_R'\\
     s_R'
\end{array}\right)
\quad
\left(\begin{array}{c}
     d_L\\
     s_L
\end{array}\right)=V_L^\dagger\left(\begin{array}{c}
     d_L'\\
     s_L'
\end{array}\right)
\]
At this point, we look at the kinetic term:
\begin{align*}
\pazocal{L}_{kin}&\supset Z_\mu\sum_f\bar{f}i\gamma^\mu fg_f\\
&\supset\left(\begin{array}{c}
    \bar{d}_{L/R}' & \bar{s}_{L/R}'
\end{array}\right)V_{L/R}
\left(\begin{array}{cc}
    g_Z^{d_{L/R}} & 0 \\
    0 & g_Z^{s_{L/R}}
\end{array}\right)(i\gamma^\mu Z_\mu)V_{L/R}^\dagger
\left(\begin{array}{c}
    d_{L/R}' \\
    s_{L/R}' 
\end{array}\right)
\end{align*}
The couplings are given by $g_f=g\cos\theta_wT^3-g'\sin\theta_wY$, hence what we obtain is:
\[
\left\{
\begin{aligned}
&g_Z^{d_R}=g_Z^{s_R}=g\cos\theta_W\cdot0-g'\sin\theta_W(-1/3)\\
&g_Z^{d_L}=g\cos\theta_W(-1/2)-g'\sin\theta_W(-1/6)\\
&g_Z^{s_L}=g\cos\theta_W\cdot0-g'\sin\theta_W(-1/3)
\end{aligned}
\right.
\]
The crucial point is that the couplings in the right sector are identical, $V_R\left(\begin{array}{cc}
    g_Z^{d_R} & 0 \\
    0 & g_Z^{s_R}
\end{array}\right)V_R^\dagger\sim\mathbb{1}g_Z^{d_R}$ and therefore we have no flavour violation. In the left sector instead we parametrize $V_L$ and find:
\[
V_L=\left(\begin{array}{cc}
    \cos\theta & \sin\theta \\
    -\sin\theta & \cos\theta
\end{array}\right)
\Rightarrow
V_L\left(\begin{array}{cc}
    g_Z^{d_L} & 0 \\
    0 & g_Z^{s_L}
\end{array}\right)V_L^\dagger=\left(\begin{array}{cc}
    g_Z^{d_L}\cos^2\theta+g_Z^{s_L}\sin^2\theta & (g_Z^{d_L}-g_Z^{s_L})\cos\theta\sin\theta \\
    (g_Z^{d_L}-g_Z^{s_L})\cos\theta\sin\theta & g_Z^{d_L}\sin^2\theta+g_Z^{s_L}\cos^2\theta
\end{array}\right)
\]
The flavour violation is given by the terms proportional to $(g_Z^{d_L}-g_Z^{s_L})$.
\begin{figure}[h]
    \centering
    \includegraphics[width=0.8\textwidth]{Images/FVZ_boson.pdf}
    \caption*{}
    \label{fig:my_label}
\end{figure}\\
Diagram (a) refers to $g_Z^{d_L}\cos^2\theta+g_Z^{s_L}\sin^2\theta$, (b) to $g_Z^{d_L}\sin^2\theta+g_Z^{s_L}\cos^2\theta$ and the flavour violating one (c) to $(g_Z^{d_L}-g_Z^{s_L})\cos\theta\sin\theta$. 
\newpage
In fact, in our toy model we have:
\begin{table}[h]
    \centering
    \begin{tabular}{c|cc}
     & G & H \\
     \hline
     Right sector & $\begin{cases}
     d_R\sim(3,1,-1/3)\\
     s_R\sim(3,1,-1/3)
     \end{cases}$ & (3,-1/3)\\
     \hline
     Left sector & $\begin{cases}
     d_L\sim(3,2,1/6)\\
     s_L\sim(3,1,-1/3)
     \end{cases}$ & (3,-1/3)\\
     \hline
    \end{tabular}
    \caption{No flavour violation in the right sector, flavour violation in the left sector.}
    \label{tab:my_label}
\end{table}\\
In the SM, there is no flavour violation because $g_Z^{d_L}=g_Z^{s_L}$. The left sector is a quark doublet (3,2,1/6) while in this toy model we have $s_L$ and $d_L$ coming from different irreps of the unbroken group.
\subsection{FCNC at loop-level}
So far, we studied everything at tree-level and we have seen how at this level in the SM there is no FCNC. Now, we move to loop level and see what happens. These processes are highly suppressed for 3 main reasons:\marginnote{\includegraphics[]{Images/GIM.jpg}\\
From \href{https://cerncourier.com/a/50-years-of-the-gim-mechanism/}{Cern Courier}: Hong-Jian He, John Ellis, John Iliopoulos, Sheldon Lee Glashow, Verónica Riquer and Luciano Maiani at a celebration of 50 years of the GIM mechanism in Shanghai.}
\begin{itemize}
    \item Loop-level contributions are usually small
    \item The structure of the $V_{CKM}$ matrix
    \item The \href{https://en.wikipedia.org/wiki/GIM_mechanism}{Glashow–Iliopoulos–Maiani (GIM) mechanism}
\end{itemize}
[SENTI \textbf{BENE} LA REGISTRAZIONE CHE NON C'HO CAPITO UN CAZZO]\\
As an example, we study the process $s\to de^+e^-$, which is represented at one-loop by the following diagrams:
\begin{figure}[h]
    \centering
    \includegraphics[width=0.8\textwidth]{Images/FCNC_loop.pdf}
    \caption*{}
    \label{fig:my_label}
\end{figure}\\
The amplitude is proportional to the $V_{CKM}$ matrix, in particular\\
$\pazocal{M}\propto\sum_{i=u,c,t}V_{is}^*f(m_i)V_{id}=\sum_{i=u,c,t}(V^\dagger)_{si}\Tilde{f}(x_i)V_{id}$ where $x_i=m_i^2/m_W^2$ and $\Tilde{f}(x_i)$ is the Inami-Lim function. When the GIM mechanism was proposed, there were only three quarks (up, down and strange). However, it required the existence of a fourth quark, the charm: the introduction of the fourth quark restored unitarity of the flavour mixing matrix, allowing to drop constant terms in $\Tilde{f}$. How? If $\Tilde{f}=$constant, then we have $\sum_{i=u,c,t}=(V^\dagger)_{si}V_{id}=\delta_{sd}=0$. This gives us a non-trivial information about the quarks running in the loop. Suppose to have only two families: the $V_{CKM}$ matrix becomes:
\[
V_c=\left(\begin{array}{cc}
    \cos\theta_c & \sin\theta_c \\
    -\sin\theta_c & \cos\theta_c
\end{array}\right)
\]
In the loop we have $(m_{u/c}/m_W)^2\ll1$ and we can expand $\Tilde{f}$ and compute the amplitude:
\[
\left\{
\begin{aligned}
&\Tilde{f}(x)=\cancelto{\text{constant}}{\Tilde{f}(0)}+\Tilde{f}'(0)x+\pazocal{O}(x^2)\\
&\pazocal{M}\propto V_{us}^*V_{ud}\frac{m_u^2}{m_W^2}+V_{cs}^*V_{cd}\frac{m_c^2}{m_W^2}=\cos\theta_c\sin\theta_c\left(\frac{m_c^2-m_u^2}{m_W^2}\right)
\end{aligned}
\right.
\]
The prediction of the charm mass was historically made studying the $K\bar{K}$ process, saturating the detector and then fitting for $m_c^2$ obtaining $m_c\sim1.5$\,GeV. In doing this, they were lucky because they managed to obtain $\pazocal{O}(1)$ numbers even ignoring the existence of the top.\marginnote{
\includegraphics[]{Images/martellone.jpg}\\
Glashow, Iliopoulos and Maiani after correctly predicting the charm mass.} By not including the top and still getting a correct result, we may ask ourselves: what is more important, the top or the charm? Let's have a look at the amplitudes:
\[
\pazocal{M}_t\propto\underset{\mathclap{\tikz \node {$\uparrow$} node [below=1ex] {\footnotesize  $\pazocal{O}(1)$};}}{\frac{m_t^2}{m_W^2}}\overset{\mathclap{\tikz \node {$\downarrow$} node [above=1ex] {\footnotesize $\sim\lambda^5$};}}{V_{td}V_{ts}^*} \quad \pazocal{M}_c\propto\underset{\mathclap{\tikz \node {$\uparrow$} node [below=1ex] {\footnotesize  $\sim(10^{-2})^2$};}}{\frac{m_c^2}{m_W^2}}\overset{\mathclap{\tikz \node {$\downarrow$} node [above=1ex] {\footnotesize $\sim\lambda$};}}{V_{cd}V_{cs}^*}
\]
This is a non-trivial fight between masses and CKM matrices: they turn out to be of similar order. Why is heavy physics contributing a lot at low energy? Why does it not decouple? In the SM, $m_t=Y_tV$: we cannot send it to large values and assuming that we are decoupling physics, we are touching the Yukawa too, i.e. increasing a coupling. 
\end{document}